% AttnRes: previously produced layer/block outputs are scored by a learned
% pseudo-query, softmax-weighted into the current layer input, then this layer
% creates a new output that can be read by later layers.
\documentclass[border=2mm]{standalone}
\usepackage{amsmath,amssymb}
% Chinese labels in figures; fonts are shared with the book.
\usepackage[UTF8,fontset=none]{ctex}
% CJK fonts for the book, the figures and the slides.
% Noto Sans CJK SC (SIL Open Font License) is the same typeface as Source Han Sans SC.
% macOS: brew install --cask font-noto-sans-cjk-sc
% Debian/Ubuntu: apt install fonts-noto-cjk
\setCJKmainfont[BoldFont={Noto Sans CJK SC Bold}]{Noto Sans CJK SC}
\setCJKsansfont[BoldFont={Noto Sans CJK SC Bold}]{Noto Sans CJK SC}
\setCJKmonofont{Noto Sans CJK SC}


% Shared visual language for the figures of the book.
% Every figure in this directory loads this file so the chapter reads as one system.
% NOTE: no blank lines are allowed inside \tikzset{...} (a blank line emits \par
% and breaks key scanning).
\NeedsTeXFormat{LaTeX2e}
\ProvidesPackage{figstyle}[2026/10/21 Figure style for Free Energy to LLMs]

\RequirePackage{tikz}
\usetikzlibrary{arrows.meta,positioning,calc,fit,shapes.geometric,shapes.misc,%
  backgrounds,matrix,decorations.pathreplacing,decorations.markings,patterns.meta}

% ---------------------------------------------------------------- palette
% Same values as book/preamble/style.tex, prefixed with F.
\definecolor{FNavy}{HTML}{14324A}
\definecolor{FTeal}{HTML}{0E7490}
\definecolor{FSky}{HTML}{E7F5F8}
\definecolor{FOrange}{HTML}{D97706}
\definecolor{FAmber}{HTML}{FFF3D6}
\definecolor{FInk}{HTML}{1F2937}
\definecolor{FSlate}{HTML}{64748B}
\definecolor{FMist}{HTML}{F5F7FA}
\definecolor{FRose}{HTML}{9F1239}
\definecolor{FRoseMist}{HTML}{FFF1F2}

% ---------------------------------------------------------------- semantics
% Teal   = probability / model / data object
% Orange = training objective, cost, or the quantity moved through the system
% Navy   = neutral structure
% Slate  = secondary label or metadata
% Rose   = failure mode, hazard, masked-out region
\tikzset{
  % Lay figures out in millimetres.  A figure designed 168mm wide is included at
  % \textwidth, so its fonts appear at exactly the size written here.
  every picture/.style={x=1mm,y=1mm,font=\small,line cap=round,line join=round},
  % ---- block styles
  nbase/.style={draw,rounded corners=1.2pt,align=center,inner sep=2.4pt,%
    minimum height=7mm,text=FInk},
  nteal/.style={nbase,draw=FTeal,fill=FSky},
  norange/.style={nbase,draw=FOrange,fill=FAmber},
  nnavy/.style={nbase,draw=FNavy,fill=white},
  nsoft/.style={nbase,draw=FSlate!55,fill=FMist},
  nrose/.style={nbase,draw=FRose,fill=FRoseMist},
  ghost/.style={draw=none,fill=none,align=center,inner sep=1pt},
  plain/.style={draw=none,fill=none,align=center,inner sep=0pt},
  % ---- arrows
  ar/.style={-{Latex[length=1.7mm,width=1.35mm]},draw=FSlate,line width=0.75pt},
  ard/.style={ar,dashed},
  arteal/.style={ar,draw=FTeal},
  arorange/.style={ar,draw=FOrange},
  arrose/.style={ar,draw=FRose},
  thickar/.style={ar,line width=1.2pt},
  % ---- text roles
  ftitle/.style={font=\large\bfseries,text=FNavy,align=center},
  ptitle/.style={font=\small\bfseries,text=FNavy,align=center},
  ptitleo/.style={font=\small\bfseries,text=FOrange,align=center},
  lbl/.style={font=\scriptsize,text=FSlate,align=center,inner sep=1pt},
  lblt/.style={lbl,text=FTeal},
  lblo/.style={lbl,text=FOrange},
  lbln/.style={lbl,text=FNavy},
  lblr/.style={lbl,text=FRose},
  sub/.style={font=\scriptsize,text=FSlate!90,align=center,inner sep=1pt},
  monob/.style={font=\ttfamily\scriptsize,text=FInk,inner sep=1pt},
  mathb/.style={font=\small,text=FInk,inner sep=1pt},
  eqn/.style={font=\scriptsize,inner sep=1.5pt,text=FInk},
  % ---- decorations
  brace/.style={decorate,decoration={brace,amplitude=2.4pt},draw=FSlate,%
    line width=0.6pt},
  bracet/.style={brace,draw=FTeal},
  braceo/.style={brace,draw=FOrange},
  axis/.style={draw=FSlate!55,line width=0.6pt},
  guide/.style={draw=FSlate!45,line width=0.5pt,dashed},
  halolabel/.style={font=\scriptsize,text=FInk,fill=white,fill opacity=0.85,%
    text opacity=1,inner sep=1.4pt,align=center},
}

% Figure title above the drawing, anchored at the figure's top-left.
\newcommand{\figtitle}[2][0]{%
  \node[ftitle,anchor=north west,outer sep=0pt] at (#1,0) {#2};}

% Rounded panel behind a subgroup of nodes.
\newcommand{\figpanel}[3][]{%
  \node[draw=FSlate!35,fill=#3,rounded corners=2pt,inner sep=2.6mm,#1] {#2};}

\endinput

\begin{document}
\begin{tikzpicture}
  \node[ptitle,anchor=west] at (0,95) {Kimi K3 / AttnRes：已有层输出汇成当前层输入};
  \node[lbln] at (31,84) {同一 token 的已有层块输出 $b_j$};
  \node[lblo] at (69,84) {当前层的 $w_l$};
  \node[lblo] at (88,72) {打分 $\to$ softmax};

  % Earlier outputs for one token are retained as separate depth entries.
  \draw[arteal] (7,15)--(7,68);
  \node[lbl,rotate=90] at (3,42) {由浅到深};
  \foreach \d in {0,...,5} {
    \pgfmathtruncatemacro{\tone}{43+mod(11*\d,38)}
    \filldraw[fill=FOrange!\tone!white,draw=white,line width=.24pt]
      (60+3*\d,75) rectangle +(3,5.5);
  }
  \draw[FOrange,line width=.7pt] (69,74)--(69,17);

  \foreach \y/\lab/\a/\strength in {20/0/0/25,34/1/1/75,48/{l\!-\!1}/{l\!-\!1}/48,62/l/l/62} {
    \ifnum\y=62
      \node[mathb,anchor=east] at (17,\y) {$b_l^{\mathrm{part}}$};
    \else
      \node[mathb,anchor=east] at (17,\y) {$b_{\lab}$};
    \fi
    \foreach \d in {0,...,5} {
      \pgfmathtruncatemacro{\tone}{38+mod(13*\d+3*\y,42)}
      \filldraw[fill=FTeal!\tone!white,draw=white,line width=.23pt]
        (20+4*\d,\y-3) rectangle +(4,6);
    }
    \draw[arteal] (45,\y)--(64,\y);
    \node[draw=FOrange,fill=FAmber,circle,minimum size=6mm,inner sep=0pt,
      font=\scriptsize,text=FOrange] at (69,\y) {$a_{\a}$};
    \draw[arorange,line width=.72pt] (74,\y)--(102,\y);
  }

  % The weights are jointly normalized over accessible outputs, then mixed.
  \draw[FOrange,line width=.7pt] (102,20)--(102,62);
  \draw[arorange] (102,41)--(109,41);
  \node[draw=FOrange,fill=FAmber,circle,minimum size=7mm,inner sep=0pt,
    font=\small,text=FOrange] at (113,41) {$\sum$};
  \draw[arteal] (117,41)--(124,41);
  \foreach \d in {0,...,5} {
    \pgfmathtruncatemacro{\tone}{43+mod(7*\d,35)}
    \filldraw[fill=FTeal!\tone!white,draw=white,line width=.24pt]
      (125+3*\d,38) rectangle +(3,6);
  }
  \node[lbln] at (134,52) {第 $l$ 层输入};
  \node[mathb] at (134,32) {$r_l=\sum_j a_jb_j$};

  \draw[arteal] (144,41)--(150,41);
  \node[nteal,minimum width=13mm,minimum height=8mm] at (158,41) {$\mathcal F_l$};
  \draw[arteal] (166,41)--(172,41);
  \foreach \d in {0,...,5} {
    \pgfmathtruncatemacro{\tone}{44+mod(9*\d,35)}
    \filldraw[fill=FTeal!\tone!white,draw=white,line width=.24pt]
      (173+3*\d,38) rectangle +(3,6);
  }
  \node[lbln] at (182,52) {第 $l$ 层输出};
  \node[lbl] at (182,32) {供后续层读取};
\end{tikzpicture}
\end{document}
